\section{Related Work}
\label{sec:related-work}

\looseness=-1
\paragraph{Self-consistency and ensemble theory.}
Self-consistency (SC) selects the most frequent answer across chain-of-thought traces \citep{wang2023sc}. Extensions address free-form generation \citep{chen2024usc}, confidence weighting \citep{taubenfeld2025cisc}, diverse prompting \citep{li2023diverse}, and reasoning-aware filtering \citep{wan2024rasc, toh2024notallvotes, kim2026reasc}. SC gains plateau early \citep{loo2025reevaluating, byerly2025selfconsistency}. The Condorcet Jury Theorem \citep{condorcet1785} guarantees majority-vote convergence under independence; ensemble theory \citep{dietterich2000ensemble, krogh1994neural} establishes that ensemble gain requires error diversity. \citet{ladha1992correlated} show correlated voters reduce effective ensemble size; our $K_\text{eff}$ adapts Fleiss' $\kappa$ from diversity metrics \citep{kuncheva2003measures, brown2005diversity} to quantify this compression (\S\ref{sec:pipeline}). \citet{kim2025correlated} demonstrate correlated errors in LLM outputs; \citet{kuai2026independent} propose cross-model independence auditing, complementary to our within-model process-$\kappa$.

\looseness=-1
\paragraph{Neuro-symbolic reasoning.}
We distinguish three paradigms: \emph{specification-generation} \citep{pan2023logiclm, olausson2023linc, ye2024satlm} and \emph{program-generation} \citep{gao2023pal, chen2023program, lyu2023faithful} formalize \emph{before} answer generation; \emph{trace-extraction} (this work) operates post-hoc, inheriting confirmation bias (\S\ref{sec:confirmation-bias}). Recent extensions include symbolic planning and verified reasoning \citep{ryu2025clover, xu2026dylos, li2025aristotle, wang2025fover, singh2026verge, luo2025logicagent, miya2025eidoku}; \citet{amonkar2025realitycheck} show formalization quality degrades with problem complexity; \citet{chan2026position} argue logical soundness alone is unreliable. Prior studies note high formalization error rates but treat them as quality problems; we show the errors are \emph{systematically biased} toward the model's answer (Observation~1; Figure~\ref{fig:causal-dag}). Our analysis reveals a deeper mechanism underlying \citet{huang2024selfcorrect}'s finding that LLMs cannot self-correct without external feedback, showing that even Z3-validated constraints carry no information beyond majority voting.

\looseness=-1
\paragraph{Distinction from CoT confirmation bias.}
Symbolic confirmation bias (SCB) is complementary to, but structurally distinct from, CoT-level confirmation bias \citep{wan2025unveiling, turpin2023language, jhaveri2026failing}. CoT bias operates at the \emph{trace level}---individual chains may not faithfully reflect the model's computation---while SCB operates at the \emph{aggregation level}: even if each trace is locally faithful, the \emph{collection} of self-extracted constraints mirrors the answer distribution.

\looseness=-1
\paragraph{Relationship to \citet{huang2024selfcorrect}.}
\citeauthor{huang2024selfcorrect} empirically demonstrate that LLMs cannot self-correct reasoning without external feedback; our work reveals a deeper mechanism underlying this finding. Where they report the phenomenon through evaluation of prompting strategies, we formalize the failure via the $M{\to}T{\to}A{\to}C$ Markov chain (Observation~1) that explains \emph{why} self-correction fails. Where they lack tools for a priori prediction, we provide three zero-cost diagnostics ($K_\text{eff}$, BR, $C_0$) that predict failure modes before deployment. Where they do not test formal verification, we systematically evaluate 27 training-free variants including Z3-validated constraints, showing that even provably sound solvers cannot escape the bottleneck.

\looseness=-1
\paragraph{Test-time compute and self-verification.}
Scaling test-time compute via additional traces \citep{snell2025scaling}, structured search \citep{yao2023tree}, or extended reasoning (DeepSeek-R1 \citep{guo2025deepseekr1}, QwQ \citep{qwen2024qwq}, OpenAI o1/o3) faces the same correlation challenge when traces share a single model's biases. Multi-agent debate \citep{liang2023encouraging, du2023improving}, self-evaluation and refinement \citep{xie2023selfevaluation, weng2023selfverification, madaan2023selfrefine, dhuliawala2023chainofverification}, iterative self-improvement \citep{shinn2023reflexion, zhou2023lats}, tool-augmented critique \citep{gou2024critic}, formal CoT verification \citep{zhu2025vericot, chen2025sets}, and multi-agent verification \citep{lifshitz2025mav} report gains, but none tests self-extracted constraints under formal solvers \citep[cf.][]{kamoi2024selfcorrect}. We find the bottleneck lies in constraint extraction, not verification: self-extracted constraints inherit the model's bias regardless of solver soundness (\S\ref{sec:confirmation-bias}). Oracle controls confirm the verification mechanism itself is sound (Appendix~\ref{app:oracle}).

\looseness=-1
\paragraph{Training-based verification.}
Trained verifiers escape the self-extraction bottleneck via external reward signals: outcome and process reward models \citep{cobbe2021gsm8k, lightman2024prm, wang2024mathshepherd}, discriminator-guided decoding \citep{khalifa2023grace}, RL-based self-verification \citep{zeng2025s2r}, dual preference optimization \citep{dupo2025}, and execution-based filtering \citep{li2022alphacode, chen2023codet}. These operate outside the pipeline of Observation~1, potentially breaking the DPI bound by violating the trace-determined assumption. The failure modes differ: SICA's bottleneck is extraction fidelity (constraints echo the model's prior), whereas PRM/ORM's is reward model generalization \citep{wang2024mathshepherd}. Our diagnostics (BR, $K_\text{eff}$) may extend to learned verification by quantifying whether learned scores introduce genuine independence.